%=======================================================================
%  CH 6 — NUMERICAL PROBLEMS & SOLUTIONS
%=======================================================================
\section*{\color{acc}Ch 6 --- Anomaly / Fraud Detection --- Numerical Problems \& Solutions}
\vspace{-4pt}{\color{acc}\rule{\textwidth}{1.1pt}}\vspace{4pt}

{\color{sub}\footnotesize This chapter is \textbf{almost entirely theory}. Across all 17 papers
there is exactly \textbf{one} numerical \yr{76 Ch}, on the base rate fallacy. It is solved below,
plus 2 practice problems covering the patterns a future numerical would most likely take.}

\lead{Methods}
\begin{itemize}
  \item \textbf{M1} Base rate fallacy: detection rate, false alarm rate, Bayesian detection rate
        \hfill \yr{76 Ch}
  \item \textbf{M2} Distance-based outlier: the DB$(r,\pi)$ test \hfill \emph{practice}
\end{itemize}

\hr

%=======================================================================
\T{M1. Base rate fallacy}

\lead{Method}
\begin{itemize}
  \item Let $I$ $=$ anomalous, $\lnot I$ $=$ normal, $A$ $=$ alarm raised.
  \item \textbf{Detection rate} $= P(A|I)$ \qquad \textbf{False alarm rate} $= P(A|\lnot I)$
  \item \textbf{Bayesian detection rate} (what an analyst actually cares about):
        \[ P(I|A) = \frac{P(A|I)\,P(I)}{P(A|I)P(I) + P(A|\lnot I)P(\lnot I)} \]
\end{itemize}
\begin{itemize}
  \item \mk{Fastest safe route}: invent a round population (10,000), build the confusion matrix,
        read every rate off it. No formula to misremember.
\end{itemize}

\creamq{Problem 1.1}{\m{3}}
\begin{asked}{2076 Chaitra \textperiodcentered{} Q11 (b) \textperiodcentered{} [3]}
If the probability that a normal object is classified as an anomaly is 0.01 and the probability
that an anomalous object is classified as anomalous is 0.99, then what is the false alarm rate
and detection rate if 99\% of the objects are normal?
\end{asked}

\lead{Given}
\begin{itemize}
  \item $P(A|\lnot I) = 0.01$ \quad (normal object raises an alarm)
  \item $P(A|I) = 0.99$ \quad (anomalous object raises an alarm)
  \item $P(\lnot I) = 0.99$, \ so $P(I) = 0.01$
\end{itemize}

\lead{Step 1: the two rates asked for, read straight off the givens}
\begin{itemize}
  \item \textbf{False alarm rate} $= P(A|\lnot I) = \mathbf{0.01}$ \ (1\%)
  \item \textbf{Detection rate} $= P(A|I) = \mathbf{0.99}$ \ (99\%)
\end{itemize}

\lead{Step 2: why the 99\% prior was given --- build the matrix on 10,000 objects}
\begin{tabularx}{0.86\textwidth}{@{}L{3.6cm}XXL{2.1cm}@{}}
\toprule
\hrow & \thead{Alarm} & \thead{No alarm} & \thead{Total} \\
Anomalous ($I$)   & $0.99\times100=\mathbf{99}$ (TP) & 1 (FN)   & 100 \\
Normal ($\lnot I$)& $0.01\times9900=\mathbf{99}$ (FP) & 9801 (TN)& 9,900 \\
Total & 198 & 9,802 & 10,000 \\
\bottomrule
\end{tabularx}

\begin{itemize}
  \item Check: detection rate $=99/100=0.99$ \checkmark \quad
        false alarm rate $=99/9900=0.01$ \checkmark
\end{itemize}

\lead{Step 3: Bayesian detection rate}
\[ P(I|A) = \frac{(0.99)(0.01)}{(0.99)(0.01) + (0.01)(0.99)}
          = \frac{0.0099}{0.0099+0.0099} = \frac{99}{198} = \mathbf{0.50} \]

\lead{Final answer}
\begin{itemize}
  \item False alarm rate $= \mathbf{0.01}$, \ detection rate $= \mathbf{0.99}$.
  \item \mk{But the Bayesian detection rate is only 0.50} --- half of every alarm raised is false,
        despite a detector that is 99\% accurate in both directions.
  \item \textbf{This is the base rate fallacy.} The normal population (9,900) is so much larger than
        the anomalous one (100) that its 1\% error rate manufactures exactly as many false alarms
        (99) as there are true detections (99).
  \item \textbf{Say this in the exam}: it is the whole point of the question and is worth more than
        the two arithmetic lines.
\end{itemize}

\begin{center}\fbox{\begin{minipage}{0.9\textwidth}\small
\mk{If the examiner wanted only the two conditional rates}, they are given directly and the
question is trivial --- which is the clue that the \textbf{99\% prior} must be used for something.
Answer both readings: state the two rates, then compute $P(I|A)$ and name the fallacy. \added
\end{minipage}}\end{center}

\creamq{Practice A --- rarer anomalies}{}
\begin{asked}[PROBLEM]{not a PYQ \textperiodcentered{} same detector, rarer anomalies \added}
The detector of Problem 1.1 is unchanged: a normal object is classified as an anomaly with
probability 0.01, and an anomalous object is classified as anomalous with probability 0.99.
Now only \textbf{1 object in 1{,}000} is anomalous instead of 1 in 100. Find the detection rate,
the false alarm rate, and the probability that an object raising an alarm really is an anomaly.
\end{asked}
{\color{sub}\footnotesize Shows how fast the Bayesian rate collapses. \added}

\begin{itemize}
  \item On 100,000 objects: 100 anomalous, 99,900 normal.
  \item TP $=0.99\times100 = 99$ \qquad FP $=0.01\times99{,}900 = 999$
  \item $P(I|A) = \dfrac{99}{99+999} = \dfrac{99}{1098} = \mathbf{0.090}$
  \item \mk{Only 9\% of alarms are real.} Detection rate and false alarm rate are unchanged at
        0.99 and 0.01 --- \textbf{only the base rate moved}.
  \item \textbf{Conclusion (Axelsson)}: to keep $P(I|A)$ usable when $P(I)$ is tiny you must drive
        the \emph{false alarm rate} down, not push the detection rate up.
\end{itemize}

\hr

%=======================================================================
\T{M2. Distance-based outlier: the DB$(r,\pi)$ test}

\lead{Method}
\begin{itemize}
  \item Object $o$ is a \textbf{DB$(r,\pi)$-outlier} if
        \[ \frac{\lVert\{o' \mid dist(o,o') \le r\}\rVert}{\lVert D\rVert} \le \pi \]
  \item In words: count how many objects (including $o$ itself, per Han \& Kamber) lie within
        distance $r$; if that count as a fraction of $|D|$ is at most $\pi$, it is an outlier.
  \item Equivalent form: with $k=\lceil \pi|D|\rceil$, $o$ is an outlier if $dist(o,o_k) > r$.
\end{itemize}

\creamq{Practice B}{}
\begin{asked}[PROBLEM]{not a PYQ \textperiodcentered{} written to fill the gap \added}
A dataset contains the points $P_1(1,1)$, $P_2(1,2)$, $P_3(2,1)$, $P_4(2,2)$ and $P_5(8,8)$.
Taking $r = 2$ and $\pi = 0.4$, and using Euclidean distance, determine which points are
DB$(2,\,0.4)$-outliers.
\end{asked}
{\color{sub}\footnotesize No PYQ has set this yet, but it is the only computable form of the
distance-based question that is asked in 6 papers, so it is worth one run-through. \added}

\lead{Step 1: distance matrix}
\begin{tabularx}{0.72\textwidth}{@{}L{1.2cm}XXXXX@{}}
\toprule
\hrow & \thead{$P_1$} & \thead{$P_2$} & \thead{$P_3$} & \thead{$P_4$} & \thead{$P_5$} \\
$P_1$ & 0 & 1 & 1 & 1.414 & 9.899 \\
$P_2$ & 1 & 0 & 1.414 & 1 & 9.220 \\
$P_3$ & 1 & 1.414 & 0 & 1 & 9.220 \\
$P_4$ & 1.414 & 1 & 1 & 0 & 8.485 \\
$P_5$ & 9.899 & 9.220 & 9.220 & 8.485 & 0 \\
\bottomrule
\end{tabularx}

\lead{Step 2: count neighbours within $r = 2$} \ (threshold $\pi|D| = 0.4\times5 = 2$ objects)
\begin{itemize}
  \item $P_1$: itself, $P_2$, $P_3$, $P_4$ $= 4$ \quad $4/5 = 0.8 > 0.4$ $\rightarrow$ normal
  \item $P_2$, $P_3$, $P_4$: 4 each by symmetry $\rightarrow$ normal
  \item $P_5$: only itself $= 1$ \quad $1/5 = 0.2 \le 0.4$ $\rightarrow$ \mk{\textbf{outlier}}
\end{itemize}

\lead{Final answer} $P_5$ is the only DB$(2,0.4)$-outlier.
\begin{itemize}
  \item \textbf{Cross-check with the $k$-NN form}: $k=\lceil 0.4\times5\rceil = 2$.
        For $P_5$, its 2nd nearest neighbour is at $8.485 > 2$ $\rightarrow$ outlier \checkmark.
        For $P_1$, its 2nd nearest is at $1 \le 2$ $\rightarrow$ normal \checkmark.
  \item \mk{Note the sensitivity to $r$}: raise $r$ to 9 and $P_5$ stops being an outlier
        altogether. Quoting the chosen $r$ and $\pi$ is part of the answer.
\end{itemize}
